\chapter{Перемежение}
\begin{center}
\textbf{Практика 3. Перемежение}
\end{center}

\noindent\textit{Задание}\\
Реализовать операцию прямого и обратного перемежения закодированного битового сообщения.

\subsubsection*{[...................................................................................................................................]}

\noindent\textit{Описание прямого перемежения}\\
Задача перемежителя – перемешать биты в кодовом слове (с выхода сверточного кодера), используя псевдослучайный закон перестановки. Внутри перемежителя должен быть описан закон перестановки (можно использовать случайный закон).


\begin{center}
\includegraphics[width=0.8\textwidth]{1_interleaving.jpg}
\end{center}
\subsubsection*{}
Для перестановки бит можно сформировать вектор, содержащий позицию бит после перестановки. Например, используя вектор бит: \(\{10010110\}\) и вектор позиций бит после перестановки \(\{85142637\}\), можно осуществить перестановку, и выходной вектор будет \(\{00110101\}\).

\begin{center}
\includegraphics[width=0.8\textwidth]{2_interleaving.jpg}
\\
\end{center}

\subsubsection*{Описание входных/выходных переменных}
\begin{tabular}{|p{0.3\textwidth}|p{0.6\textwidth}|}
\hline
\textbf{Входные переменные:} & вектор с изначальными битами \\
\hline
\textbf{Выходные переменные:} & вектор с перемешанными битами \\
\hline
\textbf{Параметры:} & нет \\
\hline
\end{tabular}\\



\noindent\textit{Описание обратного перемежителя}\\
Задача обратного перемежителя – перемешать биты в обратном порядке, используя такой же закон перестановки бит, как в перемежителе.\\
Биты на входе прямого перемежителя и на выходе обратного перемежителя должны совпадать.\\
\noindent\textit{Описание входных/выходных переменных}


\begin{tabular}{|p{0.3\textwidth}|p{0.6\textwidth}|}
\hline
\textbf{Входные переменные:} & вектор с перемешанными битами \\
\hline
\textbf{Выходные переменные:} & вектор с изначальными битами \\
\hline
\textbf{Параметры:} & нет \\
\hline
\end{tabular}


\begin{center}
\includegraphics[width=0.7\textwidth]{3_interleaving.jpg}
\\
\end{center}
\endinput